\chapter{Teoretická část}

\section{Současný stav}

Moderní mobilní operační systémy poskytují vývojářům třetích stran standardizovaná
rozhraní pro přístup ke zdravotním datům uloženým v~centralizovaném systémovém
úložišti. Na platformě Android tuto roli od verze Android~8 zastává Health
Connect~\cite{healthconnect2024}, na platformě iOS pak HealthKit~\cite{healthkit2024},
dostupný od verze iOS~8. Obě rozhraní fungují na principu zprostředkování --
aplikace nezíská přímý přístup k~datům jiných aplikací, nýbrž čte z~a zapisuje do
sdíleného úložiště, přičemž uživatel musí explicitně udělit oprávnění ke každému
datovému typu zvlášť. Tento model výrazně zjednodušuje integraci mezi aplikacemi
(sledování kroků, chytré hodinky, fitness tracker) a~eliminuje nutnost vytvářet
proprietární synchronizační protokoly.

Klíčovým modelem při návrhu mobilních aplikací závislých na vzdáleném serveru
je přístup \textit{offline-first}~\cite{offlinefirst2015}. Mobilní zařízení jsou
ze své podstaty vystavena nestabilnímu nebo zcela chybějícímu síťovému připojení,
proto robustní aplikace nepředpokládá nepřetržitou konektivitu, nýbrž primárně
pracuje s~lokálně uloženými daty a~se serverem se synchronizuje, jakmile je připojení
k~dispozici. Na platformě React Native je standardním řešením pro rychlé lokální
úložiště knihovna MMKV~\cite{mmkv2024}, která využívá nativní implementaci ve formě
sdílené paměti mapované na soubor a~dosahuje výrazně nižší latence než
\texttt{AsyncStorage}.

Komunikace mezi mobilním klientem a serverovou částí je v současných aplikacích
zpravidla realizována prostřednictvím architektury REST 
(\textit{Representational State Transfer})~\cite{fielding2000rest}. REST definuje
sadu architektonických omezení -- bezstavovost, jednotné rozhraní, vrstvený systém --
která vedou k~interoperabilním, škálovatelným a~snadno testovatelným API. Každý
zdroj je identifikován URL adresou, operace nad ním jsou vyjádřeny standardními
HTTP metodami (GET, POST, PUT, DELETE) a~data jsou přenášena zpravidla ve formátu
JSON.

Autentizace uživatelů v~moderních webových a~mobilních aplikacích se nejčastěji
řeší dvěma přístupy. Prvním jsou tokeny JWT (\textit{JSON Web Token}), které jsou
bezstavové a~snadno horizontálně škálovatelné, avšak jejich uložení na straně
klienta přináší bezpečnostní rizika, pokud jsou uchovávány v~přístupném úložišti.
Druhým přístupem jsou stavové relace spravované serverem s~identifikátorem
uloženým v~\texttt{HttpOnly} cookie, která není přístupná JavaScriptovým kódem,
čímž je eliminováno riziko krádeže tokenu skriptem třetí strany~\cite{owasp2024mobile}.

\section{Existující řešení}

Na trhu existuje řada aplikací zaměřených na sledování zdravotních metrik.
Níže jsou popsána nejrozšířenější řešení včetně jejich omezení, jež vedla
k~rozhodnutí vytvořit vlastní aplikaci.

\subsection{Google Fit}

Google Fit~\cite{googlefit2024} je platforma vyvinutá společností Google,
integrovaná přímo do operačního systému Android. Automaticky zaznamenává pohybovou
aktivitu prostřednictvím senzorů zařízení a~slouží jako centrální úložiště
zdravotních dat pro ekosystém Android. Aplikace nabízí jednoduché přehledy
a~cíle, avšak detailní sledování výživy ani spánku nativně nepodporuje --
tyto oblasti vyžadují napojení na specializované aplikace třetích stran.

\subsection{Apple Health}

Apple Health~\cite{applehealth2024} zastává v~ekosystému iOS analogickou roli
jako Google Fit na Androidu -- jde o~systémové úložiště zdravotních dat přístupné
přes rozhraní HealthKit. Aplikace agreguje záznamy z~různých zdrojů a~zobrazuje
je v~přehledném rozhraní. Platforma je dostupná výhradně na zařízeních Apple
a~funguje primárně jako datový hub; vlastní zadávání výživy nebo tréninkových
plánů není nativně podporováno.

\subsection{MyFitnessPal}

MyFitnessPal~\cite{myfitnesspal2024} je jednou z~nejrozšířenějších aplikací pro
sledování výživy na světě. Disponuje rozsáhlou databází potravin, podporou
skenování čárových kódů a~integrací s~fitness zařízeními. Sledování tréninkového
zatížení je omezeno na kalorický výdej bez možnosti podrobné analýzy silových
tréninků. Velká část funkcionalit -- pokročilé makroanalýzy, plánování, hloubkové
statistiky -- je dostupná pouze v~placené verzi Premium.

\subsection{Cronometer}

Cronometer~\cite{cronometer2024} se od konkurence odlišuje důrazem na nutriční
přesnost. Databáze potravin vychází primárně z~ověřených vědeckých zdrojů
(USDA, NCCDB) a~aplikace sleduje vedle makroživin také mikroživiny, vitamíny
a~minerály. Sledování fyzické aktivity a~spánku je v~základní verzi omezeno;
pokročilé funkce jsou vázány na placenou verzi Gold.

\subsection{Sleep Cycle}

Sleep Cycle~\cite{sleepcycle2024} je aplikace specializovaná výhradně na analýzu
kvality spánku. Ke sledování využívá mikrofon zařízení nebo propojení s~chytrými
hodinkami a~na základě detekce pohybu a~zvukových vzorů odhaduje spánkové fáze.
Nabízí chytrý budík probouzející uživatele v~nejlehčí fázi spánku. Aplikace
nenabízí žádnou integraci se sledováním výživy či pohybu a~pokročilé statistiky
jsou dostupné pouze v~rámci předplatného.

\section{Volba nástrojů a zdůvodnění}

Výběr technologií byl proveden s~ohledem na požadavky projektu,
a osobní znalostí s některými z technologií.

\subsection{React Native a Expo}

Pro vývoj mobilního částí byl zvolen framework React Native~\cite{reactnative}
v~kombinaci s~platformou Expo~\cite{expo}. React Native umožňuje sestavit jedinou
sdílenou kódovou základnu v~jazyce JavaScript/TypeScript, ze které jsou generovány nativní
aplikace pro Android i~iOS. Tato vlastnost byla klíčovým kritériem výběru --
projekt byl sice primárně navrhován pro Android, avšak podpora iOS je zachována
pro případnou budoucí expanzi bez nutnosti přepisovat aplikaci od základu.
Rozhodujícím faktorem byla také předchozí znalost ekosystému React, jež výrazně
zkrátila dobu potřebnou k~osvojení frameworku.

Jako alternativy připadaly v~úvahu framework Flutter a~nativní vývoj v~jazycích
Kotlin resp. Swift. Flutter sice nabízí srovnatelnou multiplatformnost, avšak
vyžaduje znalost jazyka Dart, který nebylo možné sdílet se serverovou částí.
Nativní vývoj by poskytl maximální přístup k~systémovým API za cenu nutnosti
udržovat dvě zcela oddělené kódové základny.

Expo bylo zvoleno jako současný standard React Native vývoje. Managed workflow
zjednodušuje konfiguraci nativních závislostí, nástroj EAS Build zajišťuje
automatizované sestavení v~rámci CI/CD pipeline na GitHub Actions a~Expo SDK
poskytuje stabilní přístup k~nativním modulům -- včetně integrace s~rozhraními
Health Connect~\cite{healthconnect2024} a~HealthKit~\cite{healthkit2024}.

\subsection{Bun a Hono}

Serverová část je postavena na rozhraní Bun~\cite{bun} s~využitím
webového frameworku Hono~\cite{hono2024}. Bun je moderní JavaScriptový runtime
s~výrazně vyšším výkonem oproti Node.js -- v~porovnání s Node.js dosahuje násobně
kratších dob startu i~vyšší propustnosti HTTP požadavků~\cite{bunvsnode2025}. Klíčovou
výhodou je vestavěná podpora jazyka TypeScript bez nutnosti transpilace, což
zjednodušuje vývojový proces při zachování plné kompatibility s~ekosystémem npm.

Hono je odlehčený webový framework optimalizovaný pro serverová prostředí.
Vyznačuje se minimálními závislostmi, typově bezpečným směrováním požadavků
a~přímočarým API. Kombinace Bun + Hono byla zvolena jako moderní alternativa
ke standardnímu stacku Node.js + Express.

\subsection{PostgreSQL}

Jako databázový systém byla zvolena databáze PostgreSQL~\cite{postgresql2024}.
Datový model aplikace LifeMeter zahrnuje vzájemně provázané entity (uživatelé,
jídla, tréninky, spánkové záznamy), jejichž vztahy jsou přirozeně vyjádřeny
relačním schématem s~cizími klíči a~referenční integritou. Relační databáze
představuje přirozenější volbu oproti dokumentovým databázím (MongoDB) nebo
embedded řešením (SQLite), která by nebyla vhodná pro serverovou část
s~souběžnými přístupy více uživatelů.

\subsection{Better Auth}

Pro správu autentizace a~uživatelských relací byla zvolena knihovna Better
Auth~\cite{betterauth2024}. Tato volba byla upřednostněna před vlastní
implementací i~před cloudovými řešeními typu Firebase Authentication. Better
Auth je navržen s~nativní podporou jazyka TypeScript a~poskytuje přímočarou
integraci s~frameworkem Hono prostřednictvím dedikovaného adaptéru.
Jako self-hosted řešení nevyžaduje závislost na infrastruktuře třetí strany,
čímž je zajištěna plná kontrola nad uživatelskými daty. Autentizace je
realizována prostřednictvím \texttt{HttpOnly} cookies se stavovými relacemi
na serveru, v~souladu s~doporučeními OWASP~\cite{owasp2024mobile}.

\subsection{Docker a Oracle Cloud}

Celý backendový systém je kontejnerizován pomocí nástroje Docker
Compose~\cite{docker2024} a~nasazen na instanci Oracle Cloud Infrastructure.
Docker Compose umožňuje deklarativní definici vícekomponentního prostředí
(API server, databáze, reverzní proxy) v~jediném konfiguračním souboru,
čímž zajišťuje reprodukovatelnost nasazení a~jednoduchou správu závislostí
mezi službami.

Oracle Cloud byl zvolen díky programu Always Free, který zahrnuje bezplatnou
instanci s~architekturou ARM s~dostatečnými výpočetními prostředky pro provoz
projektu. Na rozdíl od PaaS řešení poskytuje vlastní virtuální stroj plnou
kontrolu nad konfigurací prostředí a~správou uložených dat.